\documentclass[dvipsnames, svgnames, x11names]{beamer}
\usetheme{mytheme2024}

\title{第十届物理化学课程教学研讨会}
\subtitle{沈阳}
\author{张斌}
\institute[理工系]{河北农业大学 理工系}
\date{\the\year 年 \the\month 月}
%-----------------------------------------------------------------
\begin{document}


\maketitle


\begin{frame}
    \frametitle{开幕式}
    \begin{figure}[htpb]
        \centering
        \includegraphics[width=.66\linewidth]{IMG20240706080740}
    \end{figure}
\end{frame}


\begin{frame}
    \frametitle{开幕式}
    \begin{figure}[htpb]
        \centering
        \includegraphics[width=.66\linewidth]{IMG20240706081503}
    \end{figure}
\end{frame}


\begin{frame}
    \frametitle{孙世刚：从电化学发展谈学科交叉前言}
    \begin{figure}[htpb]
        \centering
        \includegraphics[width=.66\linewidth]{IMG20240706090219}
    \end{figure}
\end{frame}


\begin{frame}
    \frametitle{孙世刚：从电化学发展谈学科交叉前言}
    \begin{itemize}
      \item 电化学起源
      \item 化学反应与电化学反应
      \item 一次电池工作原理：Danielle cell
      \item 电化学加工制造
      \item 我国主要的电化学工业
      \item 高端制造电子电镀
      \item 集成电路的制作过程
      \item 未来30年集成电路制造的发展趋势
    \end{itemize}
    
\end{frame}


\begin{frame}
    \frametitle{孙世刚：从电化学发展谈学科交叉前言}
    
    \begin{itemize}
      \item 电化学能源
        \begin{itemize}
          \item 能量转化
          \item 能量储存
          \item 能量转化和储存
        \end{itemize}
      \item 电化学能源的发展
      \item 氢氧燃料电池
      \item 碳达峰、碳中和倒逼能源结构变革
      \item 电化学储能促进可再生能源利用
      \item 新能源汽车动力电池
    \end{itemize}
\end{frame}


\begin{frame}
    \frametitle{孙世刚：从电化学发展谈学科交叉前言}
    电化学发展及其学科前沿交叉
    \begin{itemize}
      \item 下一代动力电池性能指标
      \item 亟待解决的关键科学问题
        \begin{itemize}
          \item 材料
          \item 界面
          \item 传输
          \item 系统
        \end{itemize}
      \item 电化学机理研究
        \begin{itemize}
          \item 电化学原位谱学
          \item 原位/工况TEM研究装置设计和加工
          \item 可调谐红外自由电子激光装置
        \end{itemize}
      \end{itemize}
\end{frame}


\begin{frame}
    \frametitle{孙世刚：从电化学发展谈学科交叉前言}
    物理化学研讨会教学体会
    \begin{itemize}
      \item 在拔尖班增设讨论课，每两周一次
    \end{itemize}
\end{frame}


\begin{frame}
    \frametitle{房喻：从基础研究的重要性看教育与人才培养}
    \begin{figure}[htpb]
        \centering
        \includegraphics[width=.66\linewidth]{IMG20240706103047}
    \end{figure}
\end{frame}


\begin{frame}
    \frametitle{房喻：从基础研究的重要性看教育与人才培养}
    \begin{itemize}
      \item 从学习角度看人工智能与教育
      \item 人类文明持续发展面临前所未有的重大挑战
        \begin{itemize}
          \item 资源问题
          \item 能源问题
          \item 环境问题
        \end{itemize}
      \item 以人工智能未代表的科技能力迅猛发展，人类应对发展危机的能力不断强化
        \begin{itemize}
          \item 2016 AlphaGo
          \item 2018 DeepMind AlphaFold
          \item 2020 Robot Chemist
          \item 2023 ChatGPT
          \item 2024 Sora
        \end{itemize}
    \end{itemize}
\end{frame}


\begin{frame}
    \frametitle{房喻：从基础研究的重要性看教育与人才培养}
    \begin{itemize}
      \item AI + Chemistry
      \item 
    \end{itemize}
\end{frame}


\begin{frame}
    \frametitle{朱志昂：二学时统计力学示范课}
    \begin{itemize}
      \item 物理化学课程中为什么要讲点统计力学？
      \item 二学时要讲哪些知识点？
      \item 怎么讲？
    \end{itemize}
\end{frame}


\begin{frame}
    \frametitle{朱志昂：二学时统计力学示范课}
    为什么要讲统计力学？
    \begin{itemize}
      \item 统计力学是联系微观和宏观的桥梁
      \item 联系媒介：配分函数
      \item 统计力学可以解释热力学三大定律的本质
    \end{itemize}
\end{frame}


\begin{frame}
    \frametitle{朱志昂：二学时统计力学示范课}
    二学时要讲哪些知识点？
    \begin{itemize}
      \item 玻尔兹曼分布定律
      \item 理想气体分子各种运动形式配分函数的求算
      \item 理想气体热力学函数熵的计算
      \item 热力学三大定律的统计力学解释
    \end{itemize}
\end{frame}


\begin{frame}
    \frametitle{朱志昂：二学时统计力学示范课}
    怎么讲？
    \begin{itemize}
      \item 经典统计力学
      \item 系综
      \item 量子统计
    \end{itemize}
\end{frame}


\begin{frame}
    \frametitle{朱志昂：二学时统计力学示范课}
    \begin{itemize}
      \item 体系微观状态的描述
      \item 量子力学的描述方法
      \item 统计力学的基本定理
      \item 玻尔兹曼分布定律 
    \end{itemize}
\end{frame}


\begin{frame}
    \frametitle{朱志昂：二学时统计力学示范课}
    热力学三大定律的统计力学解释
    \begin{itemize}
      \item 内能的本质
    \end{itemize}
\end{frame}


\begin{frame}
    \frametitle{陈胜利：拔尖人才物理化学课程教学心得与思考}
    以习惯、能力和素养为目标的教学理念
    \begin{itemize}
      \item 教
      \item 学
    \end{itemize}
\end{frame}


\begin{frame}
    \frametitle{陈胜利：拔尖人才物理化学课程教学心得与思考}
    引导学生从繁杂多样的知识体系、概念及公式中跳脱出来，建立物理和化学图像，从源头和本质上理解概念，公式的内涵和外延。
    \begin{itemize}
      \item 梳理知识逻辑和主线，建立知识框架
      \item 建立知识体系的关联和统一性
      \item 建立微观机制与宏观现象的联系
      \item 基于“问题”为导向
      \item 归纳和总结知识体系的要点
      \item 基础知识与前沿研究相结合
    \end{itemize}
\end{frame}




\begin{frame}
    \frametitle{彭笑刚：适合本科教学的物理化学框架}
    \begin{itemize}
      \item 化学学科现状
    \end{itemize}
\end{frame}


\begin{frame}
    \frametitle{彭笑刚：适合本科教学的物理化学框架}
    
\end{frame}


\begin{frame}
  \frametitle{彭笑刚：适合本科教学的物理化学框架}
  
\end{frame}


\begin{frame}
  \frametitle{彭笑刚：适合本科教学的物理化学框架}
  
\end{frame}


\begin{frame}
  \frametitle{彭笑刚：适合本科教学的物理化学框架}
  
\end{frame}


\begin{frame}
  \frametitle{彭笑刚：适合本科教学的物理化学框架}
  
\end{frame}


\begin{frame}
  \frametitle{彭笑刚：适合本科教学的物理化学框架}
  
\end{frame}




































\end{document}